\section{Context and motivation}

Laboratory automation depends on operational knowledge that is difficult to preserve as a conventional dataset. A device-integration answer can depend on model, firmware, licence, communication path, and software version. A method-reproduction question can depend on source files, external libraries, calibration, resource definitions, and active runtime state. A recovery decision can depend on which physical actions completed, where material remains, and which interventions occurred. Public technical communities record these details, although the information remains distributed across threads and its public outcome can be censored when work moves to private support or local implementation.

Existing integration standards and open frameworks improve machine interoperability \cite{bar2012sila,juchli2022sila2,porr2020gateway,bromig2022manager,wierenga2023pylabrobot}. Autonomous-laboratory research adds orchestration, adaptive experiment selection, and closed-loop execution \cite{abolhasani2023rise,hung2024community,canty2023autonomy,fei2024alabos,cousty2024glas}. These systems create demand for evidence objects that retain operational applicability, provenance, uncertainty, validation stage, maintenance ownership, and correction history.

The Lab Automation Observatory addresses this need through a derived and governed research resource. It converts a purposive public evidence base into typed records, bounded metrics, audit ledgers, reusable schemas, deterministic analyses, and versioned release artifacts. The resource is designed to support independent recoding, comparative schema work, prospective operational studies, and community knowledge stewardship. It excludes a verbatim forum archive and cannot support estimates of field prevalence, product reliability, installed base, or market share.
